\section{Literature Review and Conceptual Framework}\label{sec:literature}

\subsection{Race outcomes as discrete choices}\label{sec:lit-logit}
The conditional-logit framework treats each race as a choice set in which one runner is observed to win. Runner covariates enter a latent score, and exponentiating and normalizing those scores within the field produces probabilities that sum to one. This is closely related to McFadden's conditional-logit formulation for qualitative choice and was applied directly to racetrack handicapping by Bolton and Chapman\cite{mcfadden1974,bolton1986}. The key advantage is not merely computational convenience. Race-level normalization encodes the fact that strengthening one competitor necessarily reduces the probability available to the others, whereas independent binary classifiers can produce incoherent race probabilities unless they are post-processed.

Harville's model provides a related route from winning probabilities to ordered finishes through successive without-replacement selection, while Henery and subsequent ranking research consider alternative probability structures for permutations\cite{harville1973,henery1981}. These models matter for a capstone because they clarify what a WIN model does and does not estimate. The present project estimates the first-stage winner probability only; it does not claim that a Harville construction automatically gives accurate place or exacta probabilities. Extending to ordered outcomes would require additional validation because assumptions that are reasonable for a winner market may be inappropriate for lower finishing positions.

A second strand uses flexible statistical learning to capture nonlinear relationships. Benter describes a fundamental handicapping model whose probability estimate is subsequently combined with the public estimate, emphasizing both the difficulty of rivaling public information and the need for out-of-sample assessment\cite{benter1994}. LightGBM is a gradient-boosted tree implementation based on decision-tree ensembles and histogram-oriented learning, making it suitable for nonlinear interactions in tabular data\cite{ke2017lightgbm}. Nevertheless, a conventional binary tree objective does not understand the race choice set. In this project, tree outputs are normalized within each race and compared with a conditional-logit estimator trained directly on race likelihood, allowing the empirical value of flexible learning to be assessed without hiding the structural distinction.

\subsection{Market information and racetrack efficiency}\label{sec:lit-market}
Final odds aggregate information dispersed among bettors, including information absent from a researcher's feature set. In a pari-mutuel WIN pool, prices are linked to the distribution of stakes and the pool deduction rather than being fixed independently by a bookmaker. Bolton and Chapman explicitly connect public probabilities, track take, and expected-return decisions, while Benter reports that a fundamental model must be highly sophisticated to compete with the public estimate\cite{bolton1986,benter1994}. This literature supports two benchmarks in the present study: normalized market probability for prediction and favourite/every-runner rules for economic comparison. A new model is useful only if its performance can be interpreted relative to these simple but demanding references.

The favourite--longshot literature further shows why observed returns should be examined across the odds distribution rather than summarized by a single average. Snowberg and Wolfers document systematically different returns by odds in North American data and evaluate risk-preference and misperception explanations\cite{snowbergwolfers2010}. However, findings from Western markets should not simply be transferred to Hong Kong. Jeong, Kim, and Ro review contrary or reverse-bias evidence in some Asian markets before presenting a different test on Korean racing\cite{jeong2019}. Earlier Hong Kong analysis also motivates market-specific empirical work\cite{chapman1994}. Accordingly, this report treats odds-band returns and calibration as measurements, not as an assumed favourite--longshot pattern.

Market efficiency is not a binary property that one backtest can prove or disprove. A model might approximate final probabilities closely but fail to identify prices with positive expected value; conversely, a small betting profit could arise from luck even when average probability quality is no better than the market. The chance of selecting a threshold after seeing attractive test results adds another concern. The present design therefore distinguishes a prespecified primary betting rule from exploratory threshold grids, reports the number of bets behind every return, and uses paired resampling for both predictive differences and ROI. This interpretation is more conservative than treating a positive historical balance as sufficient evidence of inefficiency.

\subsection{Proper scoring and probability calibration}\label{sec:lit-calibration}
Proper scoring rules provide a principled basis for comparing probabilistic forecasts. Gneiting and Raftery explain that such scores incentivize honest probability statements and summarize their calibration and sharpness properties\cite{gneiting2007}. Race log loss penalizes the negative log probability assigned to the observed winner, while multiclass Brier loss measures squared probability error over the entire field. Both are more informative for this project than winner-selection accuracy alone. A model that selects the same top runner can still receive a very different log loss if it assigns excessive confidence to incorrect outcomes or insufficient probability to the actual winner.

Calibration concerns agreement between stated probabilities and observed frequencies. Niculescu-Mizil and Caruana show that probability estimates from supervised learners can require post-processing and that the best calibration approach depends on the learning method and data regime\cite{niculescu2005}. Temperature scaling offers a deliberately simple one-parameter adjustment: raising probabilities to $1/T$ and renormalizing sharpens or softens the distribution while preserving within-race ranks. It cannot repair every form of local miscalibration, but its low complexity is appropriate for a single validation season. The interim analysis consequently evaluates reliability diagrams and expected calibration error alongside global scores rather than assuming that race-wise normalization alone produces calibrated forecasts.

There is also a conceptual difference between calibration and profitability. A well-calibrated model applied to every runner in a race should tend to reproduce the economics embedded in market prices, including the deduction. Betting only when model expected value exceeds a threshold introduces selection, so small calibration errors can be amplified in the selected subset. This is especially important at long odds, where a one- or two-win difference can change ROI dramatically. Reliability evidence is therefore necessary for interpreting betting output but is not sufficient; settlement, sample size, and selection uncertainty must be reported directly.

\subsection{Validation, betting decisions, and uncertainty}\label{sec:lit-validation}
Randomly dividing runner rows would allow later races and closely related entities to appear on both sides of a split, overstating generalization to a future season. Chronological evaluation instead trains on earlier observations and validates or tests on later ones, consistent with the principle underlying time-ordered cross-validation\cite{sklearn_timesplit}. The present project uses three complete training seasons, one validation season, and one final test season. Hyperparameters and temperature parameters are restricted to training and validation information. Historical features continue to accumulate across split boundaries, but only from races that would already have occurred at prediction time.

The expected-value rule follows directly from probability and decimal odds: a bet is attractive when estimated probability multiplied by gross return exceeds the stake by a sufficient margin. Yet betting evaluation introduces dependence and skew that are easy to underestimate. Races on the same meeting share weather, track condition, and market environment; long-priced selections produce infrequent large returns; and threshold searches create researcher degrees of freedom. The interim analysis therefore clusters bootstrap resampling by meeting, records maximum drawdown, and presents sensitivity as a diagnostic rather than selecting the most profitable cell. This statistical caution reflects the broader warning in racetrack literature that apparent profitability must survive genuine hold-out evaluation\cite{bolton1986,benter1994}.

Table~\ref{tab:lit-synthesis} summarizes how the literature informs specific design choices. It also identifies what the current project adds: a recent five-season Hong Kong panel, explicit same-race leakage controls, race-level and tree-based comparisons under one protocol, official-dividend settlement, and paired uncertainty estimates for both prediction and betting. The contribution is an integrated and auditable empirical workflow, not a claim to have invented conditional logit, boosted trees, calibration, or expected-value wagering.

\begin{table}[htbp]
\centering
\caption{Literature synthesis and implications for this project}\label{tab:lit-synthesis}
\small
\begin{tabularx}{\textwidth}{L{3.5cm} L{4.4cm} Y}
\toprule
\textbf{Research strand} & \textbf{Representative sources} & \textbf{Design implication} \\
\midrule
Discrete-choice and ranking models & McFadden; Harville; Bolton and Chapman; Henery\cite{mcfadden1974,harville1973,bolton1986,henery1981} & Estimate probabilities within a race and reserve ordered-finish models for a separately validated extension. \\
\addlinespace
Fundamental and combined handicapping & Benter; Chapman\cite{benter1994,chapman1994} & Measure incremental information by comparing fundamental and market-augmented specifications out of sample. \\
\addlinespace
Nonlinear tabular learning & Ke et al.\cite{ke2017lightgbm} & Include LightGBM as a flexible comparator, but normalize runner outputs before race evaluation. \\
\addlinespace
Market efficiency and odds patterns & Snowberg and Wolfers; Jeong et al.\cite{snowbergwolfers2010,jeong2019} & Test Hong Kong odds-band behaviour empirically and avoid assuming a Western favourite--longshot pattern. \\
\addlinespace
Scoring and calibration & Gneiting and Raftery; Niculescu-Mizil and Caruana\cite{gneiting2007,niculescu2005} & Use log loss, Brier loss, reliability, and validation-only temperature scaling rather than accuracy alone. \\
\addlinespace
Chronological generalization & Time-series validation documentation\cite{sklearn_timesplit} & Preserve time order in splitting and construct every historical predictor from prior races only. \\
\bottomrule
\end{tabularx}
\end{table}
